\chapter*{Glossary}
\addcontentsline{toc}{chapter}{Glossary}
\markboth{Glossary}{Glossary}

\begin{description}[
    style=nextline,
    leftmargin=2.2em,
    labelwidth=2em,
    labelsep=0.4em,
    itemsep=0.65em,
    parsep=0pt,
    topsep=0.25em
]

\item[Action] A quantity assigned to a history by integrating its Lagrangian over time. For a common mechanical system the Lagrangian is kinetic energy minus potential energy. Physical paths make the action stationary under suitable variations; they need not minimize it. See Chapter~5.

\item[AdS/CFT correspondence] Maldacena's proposed duality between certain gravity/string theories in anti-de Sitter settings and conformal field theories on their boundaries. Strongly supported by theoretical checks, not a general proof of holography for our universe. See Chapter~8.

\item[Area law] The observation that entanglement between a region and its surroundings scales with the region's boundary rather than its volume, for many physical ground states. See Chapters~4 and~8.

\item[Arrow of time] The one-way direction of time in everyday experience, tied to the increase of entropy. See Chapter~10.

\item[Asset baking] In game development, doing expensive computation such as lighting once in advance and storing the result for cheap reuse. Used in this book as an analogy for cosmic inflation. See Chapter~9.

\item[Bekenstein bound] A proposed upper limit on the entropy of suitable bounded systems of radius \(R\) and energy \(E\): \(S \le 2\pi k_B R E / \hbar c\). Its applicability depends on physical assumptions; cosmological settings require care. See Chapters~3 and~8.

\item[Bell's theorem] John Bell's 1964 proof that no theory with local, pre-set values can reproduce all the predictions of quantum mechanics. Experiments confirm the quantum predictions. See Chapter~4.

\item[Born rule] The rule that the probability of a measurement outcome equals the squared size of its quantum amplitude. See Chapter~4.

\item[Bose--Einstein condensate] A state of matter in which many bosonic atoms occupy the same quantum state and behave as a single wave. See Chapter~4.

\item[Boson] A particle, such as a photon, that can share a quantum state with identical particles. Compare \emph{fermion}. See Chapter~5.

\item[Brute-force wall] This book's name for the prohibitive cost of specified naive classical simulations. It is not a lower bound on every possible representation. See Chapter~1.

\item[Chaos] Extreme sensitivity of a system to its starting conditions, so that tiny differences grow exponentially over time. See Chapter~1.

\item[Coarse-graining] Averaging over fine details to get a simpler description at a larger scale. See Chapter~8.

\item[Cosmic event horizon] An observer-dependent boundary limiting which events can ever send that observer a signal. Its existence depends on the future expansion history; a persistent positive cosmological constant produces one. See Chapter~10.

\item[Cosmic inflation] A proposed brief period of extremely rapid expansion in the first fraction of a second after the Big Bang. Supported by observations but not yet confirmed. See Chapter~9.

\item[Cosmic microwave background] The faint, nearly uniform radiation left over from the hot early universe, now at about 2.7 kelvin. See Chapters~6 and~9.

\item[Cosmological constant] A constant energy density of empty space, the simplest explanation for dark energy. See Chapter~10.

\item[Dark energy] Whatever is causing the expansion of the universe to accelerate. See Chapter~10.

\item[Decoherence] Suppression of locally observable interference as a system becomes correlated with its environment. It requires no conscious observer and does not alone select one definite outcome. See Chapter~7.

\item[Density functional theory] A method based on the Hohenberg--Kohn theorem, which shows that ground-state properties of electrons are determined by their density in three-dimensional space. See Chapter~4.

\item[Effective field theory] A description of physics valid at a particular range of scales, with the effects of smaller scales summarized by a few parameters. See Chapter~8.

\item[Entanglement] A quantum relationship in which several particles share a single joint state that can't be split into separate states for each particle. See Chapters~4 and~8.

\item[Entropy] A measure of how many microscopic arrangements are consistent with a system's large-scale description; equivalently, the information a coarse description leaves out. See Chapter~6.

\item[Event-driven architecture] A software design in which work happens in response to events rather than through central polling. Used as an analogy for decoherence. See Chapter~7.

\item[Fermion] A particle, such as an electron, proton, or neutron, that obeys the Pauli exclusion principle. See Chapter~5.

\item[Flyweight pattern] A software technique of sharing one definition among many identical objects. Used as an analogy for identical particles. See Chapter~3.

\item[Garbage collection] Reclaiming program memory that can no longer be reached. The thermodynamic analogy concerns inaccessible information, not demonstrated reclamation of physical storage. See Chapter~6.

\item[Hawking radiation] Radiation predicted by quantum field theory on a black-hole background. Not directly confirmed for astrophysical black holes. See Chapters~8 and~10.

\item[Holographic principle] The proposal that gravitational physics admits descriptions associated with boundaries rather than independent volume elements. Strongly supported in specific duality settings; its precise application to our universe remains open. See Chapter~8.

\item[Landauer's principle] Erasing one bit of information releases at least \(k_B T \ln 2\) of heat. See Chapter~6.

\item[Lazy evaluation] Deferring a computation until its result is needed. Its use for quantum measurement here is an interpretation-dependent analogy, not an established mechanism. See Chapter~4.

\item[Light cone] The set of events that could causally influence, or be influenced by, a given event, bounded by the speed of light. See Chapter~5.

\item[Lossy compression] Compression that discards details judged unimportant, as in JPEG and MP3 files. Used as an analogy for how the universe keeps records of its past. See Chapter~10.

\item[Multigrid] A numerical method that solves a problem at several scales at once, passing corrections between coarse and fine grids. See Chapter~8.

\item[No-communication theorem] The result that entanglement can't be used to send signals: nothing done to one particle changes the measurement statistics of its distant partner. See Chapter~8.

\item[Noether's theorem] Emmy Noether's 1918 result that every continuous symmetry of the laws of physics gives a conserved quantity. See Chapter~3.

\item[Past Hypothesis] The assumption that the universe began in a state of very low entropy, which explains why entropy increases toward the future. See Chapters~9 and~10.

\item[Pauli exclusion principle] No two identical fermions can occupy the same quantum state. See Chapter~5.

\item[Planck length] About \(1.6 \times 10^{-35}\) meters, the scale at which quantum effects of gravity are expected to become important. See Chapters~2 and~3.

\item[Pointer states] The states of a system that survive interaction with the environment and therefore appear classical. See Chapter~7.

\item[Procedural generation] Creating content from compact rules and a random seed rather than storing it in full. Used as an analogy for the growth of cosmic structure. See Chapter~9.

\item[Quantum Darwinism] Wojciech Zurek's account of how information about pointer states is copied many times into the environment, letting many observers agree. See Chapter~7.

\item[Race condition] A software bug in which the outcome depends on the uncontrolled timing of separate processes. See Chapter~5.

\item[Renormalization group] Kenneth Wilson's method for tracking how a physical theory changes as you zoom out, showing that most microscopic details wash out at large scales. See Chapter~8.

\item[Scrambling] The spreading of local information throughout a system's correlations until it can't be recovered from any small part. See Chapter~10.

\item[Sharding] Dividing a computing system into separately managed pieces. The comparison with cosmic horizons is limited: horizons depend on observers and do not define fixed servers. See Chapter~10.

\item[Singularity] A failure of spacetime completeness in general relativity, often accompanied by divergent curvature in familiar solutions. Not every singularity theorem asserts infinite density at a point. See Chapters~2 and~3.

\item[Superposition] A quantum state that combines several possible outcomes, each with its own amplitude. See Chapter~4.

\item[Wavefunction] The mathematical object that assigns an amplitude to every possible configuration of a quantum system. See Chapters~1 and~4.

\end{description}
